\documentclass[11pt]{amsart}
\usepackage[margin=1in]{geometry}
\usepackage{amsmath,amssymb,amsthm,mathtools}

\newtheorem{theorem}{Theorem}
\newtheorem{proposition}{Proposition}
\newtheorem{lemma}{Lemma}
\theoremstyle{definition}
\newtheorem{remark}{Remark}
\newtheorem{setup}{Setup}

\DeclareMathOperator{\Def}{Def}
\DeclareMathOperator{\coker}{coker}
\DeclareMathOperator{\im}{im}
\DeclareMathOperator{\Aut}{Aut}
\DeclareMathOperator{\Hom}{Hom}
\DeclareMathOperator{\Ext}{Ext}
\DeclareMathOperator{\Pic}{Pic}
\DeclareMathOperator{\rank}{rank}
\DeclareMathOperator{\Spec}{Spec}
\newcommand{\cO}{\mathcal{O}}
\newcommand{\cN}{\mathcal{N}}
\newcommand{\cT}{\mathcal{T}}
\newcommand{\Om}{\Omega}
\newcommand{\PP}{\mathbb{P}}
\newcommand{\CC}{\mathbb{C}}
\newcommand{\Gm}{\mathbb{G}_m}
\newcommand{\Xt}{\widetilde{X}}

\begin{document}

\title{An obstructed point of the moduli space of stable surfaces:\\
the double cover of the cone $\PP(1,1,e)$}
\author{}
\date{\today}

\begin{abstract}
For $e\ge 5$ let $X$ be the Gorenstein stable surface arising as the degree--$2$
canonical cover of the cone $Y=\PP(1,1,e)$ embedded by $\cO_Y(e)$, with
$K_X^2=2e$, $\chi(\cO_X)=e+3$, and a single simple elliptic singularity of
degree $2e$. Granting that the reduced moduli germ at $[X]$ has dimension
$9e+19$, we prove that $[X]$ is an obstructed, non-reduced point: the Zariski
tangent space $T^1_X$ has dimension at least $10e+19$, so there are at least
$e$ first--order deformations that no arc integrates.
\end{abstract}

\maketitle

\section{Setup and statement}

\begin{setup}\label{setup}
Fix an integer $e\ge 5$. Let $Y=\PP(1,1,e)$, the cone over the rational normal
curve of degree $e$, embedded by $\cO_Y(e)$ as a surface of minimal degree $e$
in $\PP^{e+1}$, with vertex $v$ a cyclic quotient singularity of type
$\tfrac1e(1,1)$. Let
\[
X=X_{4e+4}\subset \PP(1,1,e,2e+2),\qquad w^2=f_{4e+4}(x_0,x_1,z),
\]
be the double cover branched over a general member of $|\cO_Y(4e+4)|$, so that
$\omega_X\cong\cO_X(e)$ is invertible, ample and base--point--free, and the
canonical morphism $\varphi=\varphi_{|K_X|}\colon X\to Y\subset\PP^{e+1}$ is
finite of degree $2$. Then $X$ is a normal Gorenstein stable surface with
\[
K_X^2=2e,\qquad p_g(X)=e+2,\qquad \chi(\cO_X)=e+3,
\]
smooth away from a single point $p$ lying over $v$, which is a simple elliptic
singularity of degree $2e$. Let $\rho\colon \Xt\to X$ be the minimal
resolution, with exceptional elliptic curve $Z$, $Z^2=-2e$. Concretely $\Xt$ is
the smooth double cover $\psi\colon\Xt\to \PP(F_e)=F_e$ branched along
$B\in|4C_0+(4e+4)f|$, with $Z=\psi^{-1}(C_0)$; here $C_0$ is the $(-e)$--section
and $f$ a ruling of $F_e$. One has $K_{\Xt}=\psi^*(ef)$, so $\Xt$ is minimal
properly elliptic ($\kappa=1$) over $\PP^1$.
\end{setup}

We work over $\CC$. Since $X$ is Gorenstein of index $1$, the KSB / stable
deformation functor coincides with the ordinary deformation functor $\Def_X$,
and (as $h^0(\Theta_X)=0$, see Lemma~\ref{lem:aut}) the germ of the moduli space
of stable surfaces at $[X]$ is $\Def_X$ modulo a finite group. Thus
\emph{$[X]$ is a smooth point of the moduli space if and only if $\Def_X$ is
unobstructed}, and the Zariski tangent space to the moduli germ is
$T^1_X:=\Ext^1(\Om^1_X,\cO_X)$.

\begin{theorem}\label{main}
Assume the reduced deformation functor $\Def_X^{\mathrm{red}}$ has dimension
$9e+19$. Then for every $e\ge 5$
\[
\dim T^1_X\;\ge\;10e+19\;>\;9e+19=\dim \Def_X^{\mathrm{red}}.
\]
In particular $[X]$ is an obstructed, non-reduced point of the moduli space of
stable surfaces: it carries at least $e$ obstructed first--order deformations,
and the corresponding component is everywhere non--smoothable.
\end{theorem}

\section{The singularity germ}

Let $(X,p)$ be the germ at the simple elliptic point. It is the affine cone over
$(E,M)$, where $E$ is an elliptic curve and $M\in\Pic^{d}(E)$ with $d=2e\ge 10$,
embedded by the complete linear series $|M|$ as a projectively normal elliptic
normal curve $E\subset\PP^{d-1}$ whose ideal is generated by quadrics. Let
$\cN=\cN_{E/\PP^{d-1}}$ be its normal bundle, of rank $d-2$ and degree $d^2$.
The cone carries a $\Gm$--action; we write $T^1(\nu)$ for the weight--$\nu$ graded
piece of $T^1_{(X,p)}$.

\begin{lemma}\label{lem:germ}
For the germ $(X,p)$ with $d=2e\ge10$:
\[
T^1(\nu)=0\ (\nu\le -2),\qquad \dim T^1(-1)=d=2e,\qquad \dim T^1(0)=1,\qquad
T^1(\nu)=0\ (\nu\ge 1).
\]
Hence $\dim T^1_{(X,p)}=2e+1$, graded as $T^1(-1)\oplus T^1(0)=\CC^{2e}\oplus\CC$.
\end{lemma}

\begin{proof}
For a cone over a projectively normal, quadratically presented curve the graded
pieces are
\[
T^1(\nu)=\coker\!\Big[H^0(M^{\nu+1})\otimes H^0(M)^\vee \to H^0(\cN\otimes M^{\nu})\Big].
\]
The normal bundle $\cN$ of an elliptic normal curve is semistable: $E$ is
homogeneous under the Heisenberg translations, which act transitively and
preserve $\cN$, forcing all Harder--Narasimhan slopes equal, of value
$\mu(\cN)=d^2/(d-2)$. Twisting, $\mu(\cN\otimes M^{\nu})=d(d+\nu(d-2))/(d-2)$.

\emph{$\nu\le-2$:} then $\mu(\cN\otimes M^\nu)\le \mu(\cN\otimes M^{-2})=d(4-d)/(d-2)<0$
for $d\ge5$, so by semistability $H^0(\cN\otimes M^\nu)=0$ and $T^1(\nu)=0$.

\emph{$\nu=-1$:} here $\cN\otimes M^{-1}$ has slope $2d/(d-2)>0$ and, being
semistable of positive slope on an elliptic curve, has $H^1=0$; hence
$h^0(\cN\otimes M^{-1})=\deg=2d$. The map from
$H^0(M)\otimes H^0(M)^\vee$ has image the ambient linear vector fields, of
dimension $d$ (the multiplication $H^0(\cO)\otimes H^0(M)\to H^0(M)$ being an
isomorphism onto the trivial isotypic part). Thus $\dim T^1(-1)=2d-d=d=2e$.

\emph{$\nu=0$:} the cokernel is $H^1(\Theta_E)=\CC$, the $j$--line direction;
translations in $\Pic^d(E)$ are absorbed and contribute nothing.

\emph{$\nu\ge1$:} standard vanishing for cones over quadratically normal curves.
\end{proof}

\begin{remark}
The hypothesis $d\ge10$ is essential in the vanishing $T^1(-2)=0$: the cubic cone
($d=3$, type $\widetilde E_6$) has $T^1(-2)\ne0$. The values
$T^1(-1)=d$ are confirmed by the hypersurface cases $\widetilde E_8,\widetilde
E_7,\widetilde E_6$ ($d=1,2,3$), where the graded Jacobian ring gives
$T^1(-1)=1,2,3$ respectively.
\end{remark}

\begin{proposition}[Pinkham]\label{prop:pinkham}
For $d=2e\ge10$ the reduced versal base of $(X,p)$ is the $1$--dimensional
equisingular ($j$--) line; its tangent space inside $T^1_{(X,p)}$ is $T^1(0)$.
\end{proposition}

\begin{proof}
The versal deformation may be taken $\Gm$--equivariantly. Each irreducible
component of the reduced base is $\Gm$--invariant. A component meeting a
positive--weight coordinate nontrivially would, on taking the closure of a
generic $\Gm$--orbit, yield a one--parameter deformation of positive base weight;
by Pinkham's correspondence this is a projective surface $Y'\subset\PP^{d}$ of
degree $d$ with elliptic normal hyperplane section, other than a cone --- i.e.\
an anticanonically embedded del Pezzo surface of degree $d$. These exist only
for $d\le9$. As $T^1(\nu)=0$ for $\nu\ne-1,0$ and negative weight $-1$ is
excluded by the same argument, the reduced base is the weight--$0$ line.
\end{proof}

\section{Global cohomology of the tangent sheaf}

\begin{lemma}\label{lem:aut}
$h^0(\Theta_X)=0$, and $\Aut(X)$ is finite.
\end{lemma}

\begin{proof}
Derivations lift to the minimal resolution, $\Theta_X=\rho_*\Theta_{\Xt}$, and
$\Xt$ is a minimal surface of nonnegative Kodaira dimension ($\kappa(\Xt)=1$),
hence has no nonzero global vector fields.
\end{proof}

\begin{proposition}\label{prop:h2}
$e-1\;\le\;h^2(\Theta_X)\;\le\;e.$
\end{proposition}

\begin{proof}
By Serre duality on the Gorenstein surface $X$,
$H^2(\Theta_X)^\vee=\Hom(\Theta_X,\omega_X)=H^0\big((\Om^1_X\otimes\omega_X)^{**}\big)$,
computed on $X\setminus p=\Xt\setminus Z$ by reflexivity:
\[
h^2(\Theta_X)=\lim_{m\to\infty} h^0\big(\Xt,\ \Om^1_{\Xt}\otimes\cO(K_{\Xt}+mZ)\big).
\]
Filter by pole order along $Z$. The graded quotient at level $m$ injects into
$H^0\big(Z,\ \Om^1_{\Xt}|_Z\otimes\cO(K_{\Xt}+mZ)|_Z\big)$, and
$\Om^1_{\Xt}|_Z$ is an extension of $\omega_Z=\cO_Z$ by $\cN^\vee_{Z/\Xt}$.
Using $K_{\Xt}\cdot Z=2e$, $Z^2=-2e$ and adjunction
$\cN^\vee_Z\otimes\cO(K+2Z)|_Z=\cO(K+Z)|_Z=\omega_Z=\cO_Z$, the two twisted
line bundles on $Z\cong E$ have degrees $2e(1-m)$ (the $\omega_Z$ piece) and
$2e(2-m)$ (the $\cN^\vee$ piece). For $m\ge3$ both are negative, so the
filtration stabilizes; the $m=2$ quotient is at most $h^0(\cO_Z)=1$.

The baseline $m=0$ term $h^0(\Om^1_{\Xt}\otimes\cO(K_{\Xt}))$ and the $m=1$ term
are computed through the double cover $\psi\colon\Xt\to F_e$. From
$0\to\psi^*\Om^1_{F_e}\to\Om^1_{\Xt}\to\cO_R(-R)\to0$ (with $R$ the ramification
divisor), twisted by $\psi^*(C_0+ef)=\cO(K_{\Xt}+\text{(anti-inv.)})$, the
$\iota$--invariant part is $H^0(F_e,\Om^1_{F_e}(C_0+ef))$. The relative cotangent
sequence of $F_e\to\PP^1$ gives a subsheaf $\pi^*\cO(-2f)(C_0+ef)$ with
$h^0(\PP^1,\cO(e-2))=e-1$ sections and quotient $\cO(-C_0)$ with no sections; the
anti--invariant part $\Om^1_{F_e}(-C_0-(e+2)f)$ vanishes, and the ramification
term has negative degree on the irreducible $R$. Hence the $m\le1$ levels
contribute exactly $e-1$.

Altogether $e-1\le h^2(\Theta_X)\le (e-1)+1=e$.
\end{proof}

\section{Proof of Theorem~\ref{main}}

Consider the local--to--global sequence for the cotangent complex of $X$:
\begin{equation}\label{ltg}
0\to H^1(\Theta_X)\xrightarrow{\ } T^1_X \xrightarrow{\ \mathrm{ev}\ }
H^0(\cT^1_X)\xrightarrow{\ \mathrm{ob}\ } H^2(\Theta_X),
\end{equation}
where $\cT^1_X=\mathcal{E}xt^1(\Om^1_X,\cO_X)$ is the skyscraper $T^1_{(X,p)}$ at
$p$, of dimension $2e+1$ (Lemma~\ref{lem:germ}). Write
$r=\mathrm{pr}_{-1}\circ\mathrm{ev}\colon T^1_X\to T^1(-1)=\CC^{2e}$ for the
composite with projection to the weight--$(-1)$ part.

\medskip
\emph{Step 1: $T(\Def_X^{\mathrm{red}})\subseteq\ker r$.}
Restriction of deformations to the germ gives a morphism $\Def_X\to V$, $V$ the
versal base of $(X,p)$, chosen $\Gm$--equivariantly. The composite
$\Def_X^{\mathrm{red}}\to V$ factors through $V_{\mathrm{red}}$, which by
Proposition~\ref{prop:pinkham} is the weight--$0$ line with tangent space
$T^1(0)$. Differentiating, $\mathrm{ev}\big(T(\Def_X^{\mathrm{red}})\big)\subseteq
T^1(0)=\ker\big(\mathrm{pr}_{-1}\big)$, i.e. $T(\Def_X^{\mathrm{red}})\subseteq\ker r$.

\medskip
\emph{Step 2: $r\ne0$.}
By exactness of \eqref{ltg}, $\im(\mathrm{ev})=\ker(\mathrm{ob})$, so
$\dim\ker(\mathrm{ob})=(2e+1)-\rank(\mathrm{ob})$. By Proposition~\ref{prop:h2},
$\rank(\mathrm{ob})\le h^2(\Theta_X)\le e$, whence
\[
\dim\big(\im(\mathrm{ev})\big)=\dim\ker(\mathrm{ob})\ \ge\ (2e+1)-e=e+1.
\]
A subspace of $H^0(\cT^1_X)=T^1(-1)\oplus T^1(0)$ of dimension $e+1\ge6$ cannot
lie in the line $T^1(0)$; therefore its projection to $T^1(-1)$ is nonzero, and
so $\mathrm{pr}_{-1}\circ\mathrm{ev}=r$ is nonzero. More precisely
\[
\rank r\ \ge\ \dim\im(\mathrm{ev})-\dim T^1(0)\ \ge\ (e+1)-1\ =\ e.
\]

\medskip
\emph{Step 3: the dimension estimate.}
The Zariski tangent space of a germ dominates its Krull dimension, so
$\dim T(\Def_X^{\mathrm{red}})\ge\dim\Def_X^{\mathrm{red}}=9e+19$ by hypothesis.
By rank--nullity for $r\colon T^1_X\to T^1(-1)$ and Step~1,
\[
\dim T^1_X=\dim\ker r+\rank r\ \ge\ \dim T(\Def_X^{\mathrm{red}})+\rank r
\ \ge\ (9e+19)+e\ =\ 10e+19.
\]

\medskip
\emph{Conclusion.}
Since $\dim T^1_X\ge 10e+19>9e+19=\dim\Def_X^{\mathrm{red}}$, the functor
$\Def_X$ is not reduced: there exist first--order deformations --- namely any
$v\in T^1_X$ with $r(v)\ne0$ --- lying in $T(\Def_X)$ but not in
$T(\Def_X^{\mathrm{red}})$, hence obstructed. These are exactly the
negative--weight local deformations of the elliptic cone, visible at first order
but integrating to no arc. Modulo the finite group $\Aut(X)$
(Lemma~\ref{lem:aut}), $[X]$ is therefore an obstructed, non-reduced point of the
moduli space of stable surfaces, lying on a component all of whose members carry
a non--smoothable degree--$2e$ simple elliptic singularity. \qed

\begin{remark}
If the estimates are tight --- $h^1(\Theta_X)=9e+18$, $h^2(\Theta_X)=e$, and
$\mathrm{ob}$ of full rank on a complement of $\ker(\mathrm{ob})$ --- the exact
value is $\dim T^1_X=10e+19$, a nilpotent thickening of embedding corank exactly
$e$. This $e$ matches the number of moduli by which the vertex of $\PP(1,1,e)$
deforms toward smooth scrolls, over which the honest Horikawa canonical double
covers (same invariants $(K^2,\chi)=(2e,e+3)$) reside in the Gieseker locus of
the same moduli space --- the smoothable component that these phantom directions
point toward but cannot reach.
\end{remark}

\begin{remark}[Checkpoints]
Three local inputs carry the argument and are each short verifications:
(1)\ $\dim T^1(-1)=2e$ via semistability of the normal bundle of the elliptic
normal curve of degree $2e$ (equivalently, surjectivity of
$H^0(\cO)\otimes H^0(M)\to H^0(M)$ onto the trivial part);
(2)\ $T^1(\nu)=0$ for $\nu\le-2$, again by semistability, which genuinely uses
$d\ge10$;
(3)\ the adjunction identity $\cO(K+Z)|_Z=\omega_Z=\cO_Z$ used in the $m=2$ step
of Proposition~\ref{prop:h2}.
The reduced dimension $\dim\Def_X^{\mathrm{red}}=9e+19$ is taken as input; if
$\Def_X^{\mathrm{red}}$ has a component through $[X]$ of larger dimension the
inequality only improves.
\end{remark}

\end{document}
